

\chapter{Synchronous Reference Flash ADC} \label{chap:synchronous}


%                      -
\section{Motivation for Architectural Selection}
%                      -

Benchmarking a novel converter architecture requires a reference point that is close enough structurally to make the comparison genuinely informative. The synchronous reference used in this work is a 4-bit Flash ADC with rail-to-rail comparators, originally presented in \cite{flashadc_ref} and designed in a 65-nm CMOS process.
For the purpose of this comparison, the architecture was re-implemented and simulated in a 16-nm technology node, under the same conditions applied to the proposed asynchronous converter throughout this dissertation.

The choice of this particular design was not arbitrary. Several properties directly align with those of the asynchronous architecture described in Chapter~\ref{chap:async}, and it is precisely this alignment that makes the comparison meaningful. Both converters use the Flash topology, in which the input is compared in parallel against a set of fixed decision thresholds and the result is produced as a thermometer code, so the fundamental conversion principle is shared. Both are designed for 4-bit resolution with $2^4-1 = 15$ parallel comparison elements, so the complexity of the comparison stage is the same on both sides. Both were implemented in the same 16-nm technology node and operate from the same $0.8$~V supply, eliminating technology and supply voltage as sources of asymmetry. Finally, both architectures are explicitly oriented towards low-power operation, which is the design space where the asynchronous approach is expected to have the most significant impact.

One difference between the designs must nevertheless be stated explicitly so that the comparison remains honest: the input ranges are close but not identical. The StrongARM comparators of the reference resolve inputs all the way to the rails, while the complementary follower of the asynchronous front-end covers almost the entire supply range but loses functionality within a small margin of each rail, as documented in Chapter~\ref{chap:async}. This asymmetry is accounted for by characterising both converters over their respective effective input ranges, and the energy comparison is in any case performed through resolution-normalised figures of merit, in which the effective number of bits of each converter appears explicitly.

The fundamental difference, and the one this comparison is intended to illuminate, is the presence of a global clock. The synchronous reference depends on a periodic clock to coordinate its comparators and define valid output instants, while the asynchronous architecture eliminates the clock entirely. Removing the clock distribution network is one of the primary motivations for adopting asynchronous operation, since the clock tree contributes a non-trivial fraction of total power in high-speed designs. By holding the remaining parameters constant and changing only this aspect, the comparison gives a direct, controlled view of what asynchronous operation actually costs, or saves, in power and performance.


%                      -
\section{Architecture Overview}
%                      -

A Flash ADC resolves an analogue input voltage in a single conversion step by comparing it simultaneously against $2^N - 1$ evenly spaced reference levels, where $N$ is the number of output bits. For a 4-bit converter, this means 15 comparators operate in parallel at every active clock edge, each producing a binary decision about whether the input exceeds its assigned threshold. The outputs of all 15 comparators form a thermometer code (a sequence of ones followed by zeros, with no transition other than the boundary between them), which an encoder then maps to the final 4-bit binary word.

\begin{figure}[htbp]
    \centering
    %  - Replace the placeholder box below with your actual figure  -
    \fbox{\parbox{0.80\textwidth}{%
        \centering\vspace{3.5cm}%
        \textit{[Insert: top-level block diagram of the synchronous Flash ADC]}%
        \vspace{3.5cm}}}
    %
    \caption{Top-level block diagram of the synchronous 4-bit Flash ADC \cite{flashadc_ref}.}
    \label{fig:flash_arch}
\end{figure}

The reference voltages fed to each comparator are generated by a resistor ladder connected between $V_{REF+}$ and $V_{REF-}$, which in this rail-to-rail implementation are tied directly to the supply rails $V_{DD}$ and $V_{SS}$. The ladder divides the full input range into 15 equally spaced steps of one LSB each, so the $k$-th comparator from the bottom receives a reference of

\begin{equation}
    V_{REF,k} = V_{SS} + k \cdot \frac{V_{DD} - V_{SS}}{2^N},
    \qquad k = 1, 2, \ldots, 2^N - 1.
    \label{eq:ladder}
\end{equation}

Since all comparators are driven by the same clock, every decision is made within a single clock period and the output is valid only at well-defined time instants. This is the defining characteristic of a synchronous converter and stands in direct contrast to the asynchronous operation discussed in Chapter~\ref{chap:async}.


%                      -
\section{The StrongARM Latch Comparator}
%                      -

\subsection{Circuit Description and Operating Principle}

The comparator used in this architecture is the StrongARM latch, a dynamic topology that has become something of a standard reference in both high-speed receiver design and data converter circuits \cite{razavi_strongarm}. It is widely adopted for good reasons: high regeneration speed, good noise performance and essentially zero static power dissipation, all of which matter in a Flash ADC where 15 comparators operate in parallel.

\begin{figure}[htbp]
    \centering
    %  - Replace the placeholder box below with your actual figure  -
    \fbox{\parbox{0.75\textwidth}{%
        \centering\vspace{3.5cm}%
        \textit{[Insert: StrongARM latch schematic]}%
        \vspace{3.5cm}}}
    %
    \caption{Schematic of the StrongARM latch comparator \cite{flashadc_ref}.}
    \label{fig:strongarm}
\end{figure}

The circuit operates in two alternating phases controlled by the clock signal: a reset phase and an evaluation phase.


The reset phase occurs when the clock is low, the tail transistor is turned off and a pair of reset switches pull both output nodes up to $V_{DD}$. This precharges the outputs to a known state before every comparison, so that each evaluation begins from the same initial condition regardless of the previous result. Because the tail transistor is off and no current path exists through the core of the circuit, static power dissipation is identically zero during this phase.

The evaluation phase occurs when the clock rises, the tail transistor turns on and current begins to flow through the differential input pair. The input transistors steer this current according to the voltage difference $\Delta V = V_{INP} - V_{INN}$: the side with the higher input voltage conducts more, causing that output node to discharge faster than the other. Once a sufficient imbalance develops across the internal nodes, the cross-coupled inverter pair at the output enters regeneration, a positive feedback mechanism that amplifies the initial imbalance exponentially until the outputs settle at valid logic levels \cite{razavi_strongarm}. The time needed to complete the regeneration is inversely related to the input overdrive: larger differences between the input voltage and the reference threshold resolve faster, while inputs very close to the threshold require more time. This input-dependent latency is a fundamental property of dynamic comparators and is one of the key differences exploited by the asynchronous architecture.

Since current only flows during the evaluation phase, and only briefly at that, the average power consumption is directly proportional to the clock frequency. This makes the StrongARM latch well suited to low-power converter design: reducing the sampling rate reduces the power linearly.

\subsection{Rail-to-Rail Input Range}

A key property of this comparator implementation is its ability to correctly resolve input voltages across the full supply range, from $V_{SS}$ to $V_{DD}$. In a Flash ADC, the comparators assigned to the top and bottom reference levels must operate very close to the supply rails, and a comparator that loses sensitivity or gain in those regions would introduce nonlinearity at the extremes of the transfer characteristic, degrading the integral linearity of the converter. The rail-to-rail design ensures this does not happen, maintaining consistent performance across the entire input range.

The asynchronous front-end pursues the same goal with a complementary input stage, but, as documented in Chapter~\ref{chap:async}, its effective window stops slightly short of both rails due to the overdrive lost by the input pairs near the supply limits. The benchmarking methodology described in Chapter~\ref{chap:Benchmarking} therefore evaluates each converter over its own effective input range.


%                      -
\section{Clock Generation}
%                      -

Since the architecture is synchronous, a stable clock signal must drive the comparator array and determine the conversion instants. For simulation purposes, the clock was generated using a ring oscillator, which provides a compact and realistic on-chip clock source without requiring an ideal external generator.

Rather than driving the comparators directly from the oscillator output, the signal was routed through a tapered inverter buffer chain designed to model the loading and parasitics of a realistic clock distribution network. The chain consists of a sequence of inverters with progressively increasing drive strength, with each stage sized at approximately twice the fanout of the previous one. The final stage of the chain is a $\times$24 fanout inverter, whose size reflects the capacitive load that a clock signal meets when distributed through the metal routing to the ADC core: the interconnect wires, vias and comparator gate inputs that the clock must charge and discharge on every cycle.

A $\times$24 final buffer is a reasonable and representative choice for modelling the routing from a clock source to the ADC, as it captures the kind of loading one would expect in a physically compact but realistic clock distribution. Neglecting this stage would result in a power estimate that is overly optimistic: in a real implementation, the clock tree is rarely a negligible contributor to total power, and any fair assessment of the synchronous architecture must account for it.

The complete clock path used in simulation is therefore the ring oscillator, followed by the buffer chain ($\times$2 fanout per stage), the $\times$24 output buffer and, finally, the comparator array.

The power measurement methodology applied to this reference converter, together with its simulated performance figures, is presented alongside the results of the proposed asynchronous design in Chapter~\ref{chap:Benchmarking}, so that the two converters are characterised and compared within a single framework.
